\problem{2.17}{Traveling-wave and standing-wave bases}
Euler's identities imply
$\ee^{\ii kx}=\cos kx+\ii\sin kx$ and
$\ee^{-\ii kx}=\cos kx-\ii\sin kx$. Thus
\begin{equation}
A\ee^{\ii kx}+B\ee^{-\ii kx}
=(A+B)\cos kx+\ii(A-B)\sin kx.
\end{equation}
Comparing coefficients gives
\begin{equation}
\boxed{C=A+B,\qquad D=\ii(A-B)},\qquad
\boxed{A=\frac{C-\ii D}{2},\qquad
B=\frac{C+\ii D}{2}}.
\end{equation}
The two descriptions span the same two-dimensional solution
space for $E=\hbar^2k^2/(2m)$, but expose different physics:
exponentials carry directional flux, sines and cosines
describe standing waves.

\problem{2.18}{Probability current of free waves}
For a fixed positive $k$ and energy $E=\hbar^2k^2/(2m)$,
write
\begin{equation}
\Psi(x,t)=[A\ee^{\ii kx}+B\ee^{-\ii kx}]
\ee^{-\ii Et/\hbar}.
\end{equation}
The one-dimensional current is
$J=(\hbar/m)\operatorname{Im}(\Psi^*\partial_x\Psi)$.
The cross terms are real in $\Psi^*\partial_x\Psi$, so
\begin{equation}
\boxed{J=\frac{\hbar k}{m}(|A|^2-|B|^2)}.
\end{equation}
A single right-moving wave $A\ee^{\ii kx}$ has positive
current, a single left-moving wave $B\ee^{-\ii kx}$
negative current, and equal amplitudes produce a standing
wave with zero net current. An ideal plane wave is not
normalizable; $J$ describes flux per unit normalization
of the scattering amplitude.

\problem{2.19}{Complex Fourier series and the Fourier transform}
Start with an ordinary Fourier series on $[-a,a]$,
\begin{equation}
f(x)=\sum_{n=0}^{\infty}
\left[a_n\sin\frac{n\pi x}{a}+b_n\cos\frac{n\pi x}{a}\right].
\end{equation}
Using Euler's identity and taking $a_0=0$,
\begin{equation}
f(x)=\sum_{n=-\infty}^\infty c_n\ee^{\ii n\pi x/a},
\quad c_0=b_0,\quad
c_n=\frac{b_n-\ii a_n}{2},\quad
c_{-n}=\frac{b_n+\ii a_n}{2}\quad(n>0).
\end{equation}
Orthogonality of exponentials gives
\begin{equation}
\int_{-a}^{a}\ee^{\ii(m-n)\pi x/a}\dd x
=2a\delta_{mn}
\Rightarrow
\boxed{c_n=\frac1{2a}\int_{-a}^{a}
f(x)\ee^{-\ii n\pi x/a}\dd x}.
\end{equation}
Now define $k_n=n\pi/a$, $\Delta k=\pi/a$ and
$F(k_n)=\sqrt{2/\pi}\,a c_n$.
Then
\begin{equation}
f(x)=\frac1{\sqrt{2\pi}}\sum_{n=-\infty}^\infty
F(k_n)\ee^{\ii k_nx}\Delta k,\qquad
F(k_n)=\frac1{\sqrt{2\pi}}\int_{-a}^a
f(x)\ee^{-\ii k_nx}\dd x.
\end{equation}
In the limit $a\to\infty$, $\Delta k\to0$,
and for appropriate integrable functions these sums
and truncated integrals approach the Fourier transform pair
\begin{equation}
\boxed{F(k)=\frac1{\sqrt{2\pi}}\int_{\mathbb R}
f(x)\ee^{-\ii kx}\dd x,\qquad
f(x)=\frac1{\sqrt{2\pi}}\int_{\mathbb R}
F(k)\ee^{\ii kx}\dd k}.
\end{equation}
\qcheck{The finite-interval argument motivates the transform
rather than giving a rigorous existence theorem for all functions.
The Plancherel identity additionally preserves the $L^2$ norm.}

\problem{2.20}{Free evolution from a two-sided exponential}
For $\Psi(x,0)=A\ee^{-a|x|}$ with $a>0$,
\begin{equation}
1=2|A|^2\int_0^\infty\ee^{-2ax}\dd x
=\frac{|A|^2}{a}
\Rightarrow\quad \boxed{A=\sqrt a}.
\end{equation}
The Fourier coefficient is
\begin{align}
\phi(k)&=\frac{\sqrt a}{\sqrt{2\pi}}
\int_{-\infty}^\infty\ee^{-a|x|}
\ee^{-\ii kx}\dd x\\
&=\sqrt{\frac a{2\pi}}
\left[\frac1{a+\ii k}+\frac1{a-\ii k}\right]
=\boxed{\sqrt{\frac2\pi}\frac{a^{3/2}}{a^2+k^2}}.
\end{align}
Each momentum component evolves with frequency
$\omega_k=\hbar k^2/(2m)$. Therefore
\begin{equation}
\boxed{
\Psi(x,t)=\frac{a^{3/2}}{\pi}
\int_{-\infty}^\infty
\frac{\ee^{\ii kx-\ii\hbar k^2t/(2m)}}{a^2+k^2}\dd k}.
\end{equation}
When $a$ is large, the original state has width
$\sigma_x\sim1/a$ and a broad momentum distribution,
so it spreads rapidly; when $a$ is small the
wavefunction is very broad and its typical momentum
scale $\hbar a$ is small.
Neither extreme is a normalizable delta function
or uniform plane wave at the literal limiting parameter.

\problem{2.21}{Gaussian wave-packet spreading}
Take $\Psi(x,0)=A\ee^{-ax^2}$ with $a>0$.
Normalization gives $A=(2a/\pi)^{1/4}$.
Its Fourier transform is
\begin{equation}
\phi(k)=\frac{A}{\sqrt{2\pi}}
\int\ee^{-ax^2-\ii kx}\dd x
=\frac{A}{\sqrt{2a}}\ee^{-k^2/(4a)}.
\end{equation}
Completing the square in the inverse transform after
multiplication by $\ee^{-\ii\hbar k^2t/(2m)}$ yields
\begin{equation}
\boxed{\Psi(x,t)=
\left(\frac{2a}{\pi}\right)^{1/4}
\frac1{\sqrt{\gamma(t)}}
\exp\left[-\frac{ax^2}{\gamma(t)}\right]},
\qquad \gamma=1+\frac{2\ii\hbar at}{m}.
\end{equation}
Define $s=2\hbar at/m$ and
$w=a/(1+s^2)$. Since
$\operatorname{Re}(1/\gamma)=1/(1+s^2)$,
\begin{equation}
\boxed{|\Psi|^2=\sqrt{\frac{2w}{\pi}}
\ee^{-2wx^2}}.
\end{equation}
The state remains Gaussian but its width grows.
Parity gives $\avg{x}=\avg p=0$.
The second moments and uncertainties are
\begin{align}
\avg{x^2}&=\frac{1+s^2}{4a},&
\avg{p^2}&=a\hbar^2,\\
\boxed{\sigma_x}&=\boxed{\frac{\sqrt{1+s^2}}{2\sqrt a}},&
\boxed{\sigma_p}&=\boxed{\hbar\sqrt a},\\
\boxed{\sigma_x\sigma_p}&=\boxed{\frac{\hbar}{2}\sqrt{1+s^2}}.
\end{align}
The momentum uncertainty remains constant, whereas
position uncertainty increases due to free dispersion.
The lower bound is saturated at $t=0$ alone.
\qcheck{A wavepacket's density can evolve even though each
Fourier mode has a fixed energy; interference between
different energies generates the spreading.}
